+\chapter{Herramientas}

Este capítulo cuenta qué librerías, modelos preentrenados, lenguajes y
recursos de cómputo se usaron. Está partido en cinco bloques: lenguajes y
entorno, librerías de procesamiento de fMRI, modelos generativos, librerías de
evaluación y arquitecturas propias.

\section{Lenguajes y entorno}

\subsection{Python 3.12}

Python 3.12 es el lenguaje principal porque su ecosistema científico está bien maduro (NumPy, 
PyTorch, scikit-learn) y porque encaja con las versiones recientes de \texttt{diffusers} y 
\texttt{transformers}. El entorno se maneja con \texttt{venv} y un \texttt{requirements\_py312.
txt} fijado por hash para que sea reproducible.

\subsection{LaTeX}

El documento se compila con \texttt{pdfLaTeX} sobre la clase \texttt{MastersDoctoralThesis.cls} 
adaptada al formato UNI-FC.

\section{Procesamiento de fMRI}

\subsection{nibabel y nilearn}

\texttt{nibabel} se encarga de leer y escribir volúmenes de neuroimagen (NIfTI, GIfTI). 
\texttt{nilearn} pone arriba utilidades de más alto nivel para enmascarar, sacar series 
temporales y manejar atlas, que son justo lo que se necesita para aislar la máscara cortical 
visual de BOLD5000.

\subsection{h5py y scikit-learn}

Los betas ya preprocesados quedan guardados en HDF5 (\texttt{h5py}). Los dos adapters (Ridge 
Lineal y Ridge Estocástico) se ajustan con \texttt{sklearn.linear\_model.Ridge}, y las 
particiones se arman con \texttt{sklearn.model\_selection}.

\subsection{PyTorch 2.x, AMP bf16 y \texttt{torch.compile}}

PyTorch 2.x sostiene la inferencia del adapter en GPU, la integración con Stable Diffusion y 
la inyección de ruido del Ridge Estocástico. El entrenamiento va con \texttt{torch.cuda.amp} 
en bf16~\cite{kobayashi2017bf16,micikevicius2018amp} ---mismo rango dinámico que fp32 y sin 
NaN en las restas centradas--- y, cuando compensa, con \texttt{torch.compile} para fusionar 
kernels en el forward.

\subsection{Optimizador y validación}

Cuando hace falta optimizar por gradiente (por ejemplo, para calibrar el $\sigma$ del Ridge 
Estocástico), se usa \texttt{torch.optim.AdamW}~\cite{loshchilov2019adamw} con weight 
decay desacoplado. La selección de configuración se hace por pairwise accuracy sobre validación
y no por la pérdida agregada: así la métrica de validación
queda alineada con el objetivo final del adapter (discriminar) y no se premian
configuraciones que bajan MSE a costa de la geometría angular.

\subsection{Ficheros .mat (BOLD5000 ROIs)}

Se usan los \texttt{CSI\{1..4\}\_ROIs\_TR34.mat} que los autores de BOLD5000 dejaron listos. 
No correr un GLM LSS local ahorra entre tres y cuatro días de cómputo y saca una fuente de 
variabilidad que no es lo que esta tesis quiere estudiar.

\section{Modelos generativos}

\subsection{diffusers y Stable Diffusion 2.1 unCLIP}

\texttt{diffusers} (Hugging Face) expone \texttt{StableUnCLIPImg2ImgPipeline}, que acepta 
condicionamiento directo por \texttt{image\_embeds}. El scheduler \texttt{DPMSolverMultistepScheduler}
~\cite{lu2022dpmsolver} con 75 pasos, CFG~8.0 y noise level en 0, para no añadirle más 
varianza al embedding del adapter.

\subsection{transformers y CLIP ViT-L/14}

\texttt{transformers} trae \texttt{CLIPModel} y \texttt{CLIPProcessor}. Se usan para 
precalcular los embeddings de los estímulos (que son el target del adapter) y para medir 
similitud coseno CLIP entre la reconstrucción y el estímulo.

\subsection{xformers}

\texttt{xformers} añade atención memory-efficient al UNet de SD y con ello baja el pico de 
VRAM de unos 9~GB a 5--6~GB en bf16. En la práctica, esa diferencia es la que decide si la 
inferencia entra holgada o no en la RTX~4070~Ti (12~GB).

\section{Librerías de evaluación}

\texttt{lpips} para distancia perceptual, \texttt{scikit-image} para SSIM y PSNR, \texttt{
pytorch-msssim} para variantes multi-escala, y rutinas propias en \texttt{evaluation.py} para 
PixCorr e identificación pareada 2-vías.

\section{Arquitecturas y herramientas propias}

\subsection{Adapter base: Ridge Lineal --- \texttt{phase2/adapter\_ridge.py}}

Es la implementación de la línea base. Resuelve en cerrado la ecuación normal regularizada y 
expone \texttt{evaluate(...)} con $R^{2}_{\mathrm{macro}}$, $\cos$ medio y MSE. Sirve para 
diagnosticar de forma cuantitativa el régimen de fallo del mapeo lineal puro que está 
documentado en el Cap.~6.

\subsection{Adapter principal: Ridge Estocástico --- \texttt{phase2/adapter\_ridge\_stoch.py}}

Toma el predictor del Ridge Lineal y le suma una perturbación Gaussiana calibrada 
$\xi \sim \mathcal{N}(\mathbf{0},\sigma^{2}\mathbf{I})$, con $\sigma$ ajustado en validación 
por pairwise accuracy; después renormaliza el embedding sobre la hiperesfera unidad y recién 
ahí lo pasa al pipeline SD~2.1 unCLIP. La idea viene de MindEye y MindEye2, donde se ve que 
llevar el embedding a zonas más densas del manifold CLIP ayuda al decodificador. Yo lo consigo 
con un mecanismo más barato, sin entrenar un contrastivo completo.

\subsection{Entrenador y calibrador \texttt{phase2/train\_ridge.py}}

Coordina el barrido de $\lambda$ del Ridge Lineal por validación cruzada de 5 folds y el de 
$\sigma$ para el Ridge Estocástico. Funciona con AMP bf16, semilla fija (42) y reanudación 
desde checkpoint.

\subsection{Pipeline reproducible \texttt{phase2/visual\_evaluator.py}}

Encadena todo el flujo: lee el HDF5, aplica el adapter (Ridge Lineal o Ridge Estocástico), 
pasa por el pipeline SD~2.1 unCLIP y exporta la reconstrucción junto con las figuras compare 
(par estímulo--reconstrucción) y grid por sujeto.

\subsection{Parser BIDS \texttt{events\_to\_stim\_order.py}}

Pasa los \texttt{events.tsv} (estándar BIDS de OpenNeuro) al orden canónico de presentación. 
Mantiene la reproducibilidad bit a bit.

\section{Entorno de cómputo}

\subsection{Hardware}

Para escribir código y hacer \texttt{git push} me apoyo en una NVIDIA RTX~2070 (8~GB). Lo 
pesado ---entrenamiento, inferencia y evaluación--- va en una NVIDIA RTX~4070~Ti (12~GB) con 
CUDA 12.1, PyTorch 2.x, Python 3.12, AMP bf16 y \texttt{xformers}. La regla del flujo de 
trabajo no tiene matices: en la máquina local no se entrena.

\subsection{Almacenamiento y reproducibilidad}

Repositorio Git con tags por fase. Datos en \texttt{data\_bold5000/}, estímulos en \texttt{
BOLD5000\_Stimuli/}, semilla 42 fija en todos los puntos donde hay aleatoriedad (incluyendo 
el ruido $\xi$ del Ridge Estocástico durante la evaluación).

\afterpage{\blankpage}

